
Под производными финансовыми инструментами (деривативами) будем понимать
контракты, стоимость которых зависит от значений во времени
одного или нескольких базовых активов (цен акций, курсов валют,
процентных ставок и т.\,д.). К типичным примерам деривативов относятся
форвардные и фьючерсные контракты, свопы, опционы. В настоящей работе, если не оговорено иного, мы под
деривативом будем иметь в виду именно опционные контракты различных типов.

Общая задача прайсинга деривативов заключается в следующем: по заданной модели динамики
рынка и будущим денежным потокам контракта необходимо определить его
справедливую цену в текущий момент времени так, чтобы не возникало
арбитражных возможностей. 

Фиксируем конечный горизонт времени $T>0$. Для описания случайности и накопления информации во времени вводится
фильтрованное вероятностное пространство
\[
(\Omega, \mathcal F, \{\mathcal F_t\}_{t\in[0,T]}, \mathbb Q).
\]

Множество исходов $\Omega$ можно интерпретировать как множество всех
возможных сценариев развития рынка на отрезке $[0,T]$. Элемент
$\omega\in\Omega$ задаёт один возможный сценарий (траекторию) изменения
цен и других рыночных величин во времени.

Фильтрация $\{\mathcal F_t\}_{t\in[0,T]}$ --- это семейство
$\sigma$-алгебр на $\Omega$, удовлетворяющее условиям
\[
\mathcal F_s \subseteq \mathcal F_t \subseteq \mathcal F
\quad\text{для } 0 \le s \le t \le T
\]
то есть информация не убывает со временем. Каждый элемент
$\mathcal F_t$ представляет собой множество событий, определимых на
основе наблюдений и доступной информации к моменту времени $t$.
Будем предполагать, что все рассматриваемые ниже стохастические
процессы адаптированы к фильтрации $\{\mathcal F_t\}$.

Риск-нейтральная мера $\mathbb Q$ --- это вероятностная мера на
$(\Omega,\mathcal F)$, под которой дисконтированные цены всех
торгуемых активов являются мартингалами. Согласно первой фундаментальной
теореме финансовой математики, отсутствие арбитража эквивалентно
существованию эквивалентной (по отношению к реальной мере $\mathbb P$)
риск-нейтральной меры $\mathbb Q$, при которой дисконтированные цены
активов являются мартингалами.

Во многих стандартных моделях дополнительно предполагается полнота
рынка. В силу второй фундаментальной теоремы финансовой математики,
полнота рынка эквивалентна единственности эквивалентной
риск-нейтральной меры. Именно под риск-нейтральной мерой $\mathbb Q$
удобно формулировать задачу прайсинга: справедливая цена дериватива
выражается через математическое ожидание его дисконтированных будущих
выплат.

Для конкретизации модели введём безрисковый и рисковые активы.
Безрисковый актив (банковский счёт) задаётся процессом
\[
B_t = \exp\!\left(\int_0^t r_u\,du\right), \qquad t\in[0,T],
\]
где $(r_t)_{t\in[0,T]}$ --- адаптированный к $\{\mathcal F_t\}$ процесс
краткосрочной безрисковой процентной ставки, и $B_0 = 1$.

Пусть $S_t = (S_t^1,\dots,S_t^d)^\top$ --- вектор цен $d$ рисковых
активов. Предполагается, что $S_t$ является $\{\mathcal F_t\}$-адаптированным
процессом, а дисконтированные цены
\[
\tilde S_t^i = \frac{S_t^i}{B_t}, \qquad i=1,\dots,d,
\]
являются мартингалами под риск-нейтральной мерой $\mathbb Q$.

Общая задача прайсинга европейского дериватива формулируется следующим
образом. Пусть $H$ --- неотрицательная $\mathcal F_T$-измеримая
случайная величина, представляющая совокупную выплату по контракту в
момент времени $T$. Тогда справедливая цена в момент времени $t$
определяется как
\[
V_t = B_t\, \mathbb E_{\mathbb Q}\!\left[
\left.\frac{H}{B_T}\right|\,\mathcal F_t\right],
\qquad t\in[0,T],
\]
а в начальный момент времени
\[
V_0 = \mathbb E_{\mathbb Q}\!\left[
e^{-\int_0^T r_u\,du}\, H \right].
\]

Перейдём теперь к американским опционам. Американский опцион допускает
исполнение в произвольный момент времени $\tau$ на отрезке $[0,T]$.
Момент исполнения должен быть временем остановки относительно
фильтрации $\{\mathcal F_t\}$, то есть случайной величиной
$\tau:\Omega\to[0,T]$, для которой множество $\{\tau\le t\}$ принадлежит
$\mathcal F_t$ для каждого $t\in[0,T]$. Обозначим через $\mathcal T$
множество всех таких моментов остановок.

Выплата по американскому опциону при исполнении в момент $\tau$
определяется функцией $h:[0,T]\times\mathbb R^d\to\mathbb R_+$:
\[
H_\tau = h(\tau, S_\tau).
\]
Классическими примерами являются опционы на один рисковый актив
($d=1$):
\begin{itemize}
  \item американский колл с ценой исполнения (страйком) $K>0$:
  \[
  h(t, S_t) = (S_t - K)^+;
  \]
  \item американский пут с ценой исполнения (страйком) $K>0$:
  \[
  h(t, S_t) = (K - S_t)^+.
  \]
\end{itemize}


Более экзотические примеры включают, например, опционы на несколько
базовых активов, когда выплата в момент исполнения зависит от вектора
цен $(S_\tau^1,\dots,S_\tau^d)$ в этот момент, и деривативы, выплаты
которых зависят не только от значения цены в момент исполнения, но и от
предшествующей траектории цены базового актива.
Однако в основе
формальной постановки остаётся та же схема: для каждого допустимого
момента исполнения $\tau\in\mathcal T$ задаётся случайная величина
$H_\tau = h(\tau, S_\tau)$, интерпретируемая как выплата при исполнении
в момент $\tau$.


Формальная постановка задачи прайсинга американского опциона в момент
времени $0$ имеет вид
\[
V_0 = \sup_{\tau\in\mathcal T}
\mathbb E_{\mathbb Q}\!\left[
e^{-\int_0^\tau r_u\,du}\, h(\tau,S_\tau)
\right],
\]
где супремум берётся по всем моментам остановки $\tau$, допускаемым
условиями контракта. Таким образом, задача прайсинга американского
опциона сводится к задаче оптимальной остановки стохастического
процесса $(S_t)_{t\in[0,T]}$ на фильтрованном вероятностном
пространстве $(\Omega,\mathcal F,\{\mathcal F_t\},\mathbb Q)$.
